% Options for packages loaded elsewhere
\PassOptionsToPackage{unicode}{hyperref}
\PassOptionsToPackage{hyphens}{url}
\PassOptionsToPackage{dvipsnames,svgnames,x11names}{xcolor}
\documentclass[
  10pt,
  fontset=fandol]{ctexart}
\usepackage{xcolor}
\usepackage[margin=16mm]{geometry}
\usepackage{amsmath,amssymb}
\setcounter{secnumdepth}{-\maxdimen} % remove section numbering
\usepackage{iftex}
\ifPDFTeX
  \usepackage[T1]{fontenc}
  \usepackage[utf8]{inputenc}
  \usepackage{textcomp} % provide euro and other symbols
\else % if luatex or xetex
  \usepackage{unicode-math} % this also loads fontspec
  \defaultfontfeatures{Scale=MatchLowercase}
  \defaultfontfeatures[\rmfamily]{Ligatures=TeX,Scale=1}
\fi
\usepackage{lmodern}
\ifPDFTeX\else
  % xetex/luatex font selection
\fi
% Use upquote if available, for straight quotes in verbatim environments
\IfFileExists{upquote.sty}{\usepackage{upquote}}{}
\IfFileExists{microtype.sty}{% use microtype if available
  \usepackage[]{microtype}
  \UseMicrotypeSet[protrusion]{basicmath} % disable protrusion for tt fonts
}{}
\makeatletter
\@ifundefined{KOMAClassName}{% if non-KOMA class
  \IfFileExists{parskip.sty}{%
    \usepackage{parskip}
  }{% else
    \setlength{\parindent}{0pt}
    \setlength{\parskip}{6pt plus 2pt minus 1pt}}
}{% if KOMA class
  \KOMAoptions{parskip=half}}
\makeatother
\setlength{\emergencystretch}{3em} % prevent overfull lines
\providecommand{\tightlist}{%
  \setlength{\itemsep}{0pt}\setlength{\parskip}{0pt}}
\usepackage{amsmath,amssymb,mathtools,needspace}
\xeCJKDeclareCharClass{CJK}{"2032,"0400->"04FF,"2160->"216B,"2460->"2473,"25A0,"25C7,"25CB,"25CF}
\IfFontExistsTF{Arial Unicode MS}{\xeCJKsetup{AutoFallBack=true}\setCJKfallbackfamilyfont{\CJKrmdefault}{Arial Unicode MS}}{}
\setcounter{secnumdepth}{0}
\setlength{\emergencystretch}{2em}
\allowdisplaybreaks
\usepackage{bookmark}
\IfFileExists{xurl.sty}{\usepackage{xurl}}{} % add URL line breaks if available
\urlstyle{same}
\hypersetup{
  pdftitle={插值型的求积公式},
  colorlinks=true,
  linkcolor={Maroon},
  filecolor={Maroon},
  citecolor={Blue},
  urlcolor={Blue},
  pdfcreator={LaTeX via pandoc}}

\title{插值型的求积公式}
\author{}
\date{}

\begin{document}
\maketitle

\subsubsection{4.1.3 插值型的求积公式}\label{ux63d2ux503cux578bux7684ux6c42ux79efux516cux5f0f}

设给定一组节点

\[
a\leq x_0<x_1<x_2<\cdots<x_n\leq b
\]

且已知函数 \(f(x)\) 在这些节点上的值，作插值函数 \(L_n(x)\)（参见式（2.2.11））。由于代数多项式 \(L_n(x)\) 的原函数是容易求出的，取 \(I_n=\int_a^b L_n(x)\,\mathrm{d}x\) 作为积分 \(I=\int_a^b f(x)\,\mathrm{d}x\) 的近似值，这样构造出的求积公式

\[
I_n=\sum_{k=0}^n A_kf(x_k)
\tag{4.1.5}
\]

是插值型的，其中求积系数 \(A_k\) 通过插值基函数 \(l_k(x)\) 积分

\[
A_k=\int_a^b l_k(x)\,\mathrm{d}x
\tag{4.1.6}
\]

得出。

由插值余项定理（定理 2.2）即知，对于插值型的求积公式（4.1.5），其余项

\[
R[f]=I-I_n=\int_a^b\frac{f^{n+1}(\xi)}{(n+1)!}w(x)\,\mathrm{d}x,
\tag{4.1.7}
\]

其中 \(\xi\) 与变量 \(x\) 有关，\(w(x)=(x-x_0)(x-x_1)\cdots(x-x_n)\)。

如果求积公式（4.1.5）是插值型的，按式（4.1.7），对于次数不大于 \(n\) 的多项式 \(f(x)\)，其余项 \(R[f]\) 等于零，因而这时求积公式至少具有 \(n\) 次代数精度。

反之，如果求积公式（4.1.5）至少具有 \(n\) 次代数精度，则它必定是插值型的。事实上，这时式（4.1.5）对于插值基函数 \(l_k(x)\) 应准确成立，即有

\[
\int_a^b l_k(x)\,\mathrm{d}x=\sum_{j=0}^n A_jl_k(x_j).
\]

注意到 \(l_k(x_j)=\delta_{kj}\)，上式右端实际上即等于 \(A_k\)，因而式（4.1.6）成立。

综上所述，得到如下结论。

\textbf{定理 4.1} 形如式（4.1.5）的求积公式至少有 \(n\) 次代数精度的充分必要条件是，它是插值型的。

\begin{quote}
校注（agent 补充）：式（4.1.7）原扫描将导数阶数印作 \(f^{n+1}(\xi)\)，未带括号；此处按原式保留。由其引用的插值余项定理及上下文，这里表示 \(n+1\) 阶导数。
\end{quote}

\end{document}
